\section{Protegendo as Crianças: O Filtro AdGuard \faChild}

Como garantir que seu filho não acesse sites adultos, de violência extrema ou de apostas no celular? Nós vamos colocar um filtro direto no \enquote{cano} de internet do aparelho.

Diferente de aplicativos espiões pesados, essa é uma configuração \enquote{escondida} no próprio sistema do seu celular que limpa a internet antes mesmo de ela aparecer na tela.

\subsection*{Como configurar no Android:}
\begin{enumerate}
	\item Abra as \textbf{Configurações} do celular.
	\item Procure por \textbf{Rede e Internet} (ou use a Lupa \faSearch\ e digite ``DNS'').
	\item Clique na opção \textbf{DNS Privado}.
	\item Marque a opção \textit{\enquote{Nome do host do provedor de DNS}}.
	\item Digite exatamente este código (tudo minúsculo): 
	
	\begin{center}
		\textbf{\textcolor{AdguardGreen}{family.adguard-dns.com}}
	\end{center}
	
	\item Clique em \textbf{Salvar}.
\end{enumerate}

\begin{tcolorbox}[colback=NordBackground, colframe=AdguardGreen, title=\faCheckCircle\ Vantagem Extra]
	\textcolor{white}{Além de proteger as crianças contra sites perigosos, este filtro também bloqueia boa parte das propagandas em jogos e aplicativos gratuitos!}
\end{tcolorbox}